% ===== 原文 pp.135--138 =====
% [原文 p.135]
通常只剩下如下形式的体作用量
\begin{equation}
  S_{\mathrm{bulk}} = \sum_{p,i,j}\int_{\mathrm{AdS}_{d+1}} N_{ij,p}\, a^i_{p+1} \wedge \dd b^j_{d-p-1} + \mathcal{A}\left(\{a^i_{p+1}\}\right)\,, \tag{6.30}
\end{equation}
其中 $a$、$b$ 是体规范场，它们是时空上的微分形式，其次数由下标决定，$N_{ij,p} \in \Z$，而 $\mathcal{A}$ 是 $a$ 型规范场的函数。

设 $N_{ij,p}$ 对每个 $p$ 都定义非退化矩阵。那么 $S_{\mathrm{bulk}}$ 就是形如 (4.258) 的离散（高阶形式）规范理论借助连续场给出的描述。并且断言：$S_{\mathrm{bulk}}$ 就是与全息对偶的 $d$ 维 QFT 的某个（高阶形式）对称性相对应的 SymTFT。

% [Tutor Note: (6.30) 里 $N_{ij,p}\in\Z$ 与「$N_{ij,p}$ 非退化」两条合起来才是离散规范理论的判据：
%  整数性来自 $\U(1)$ 规范场的量子化（体作用量须在闭流形上模 $2\pi$ 定值），非退化保证
%  对易关系给出有限群而非连续群；这与 §4.5 中 BF 型 SymTFT 的写法 (4.258) 逐项对应。
%  若某个 $p$ 上 $N_{ij,p}$ 退化，则该扇区只剩连续规范场，不再是离散高阶形式对称性。]

\para{来自全息的对称性边界条件。}
正如 §4.5 讨论的，$d$ 维 QFT 的各种规范框架或全局形式，是通过选取 SymTFT 规范场的边界条件得到的。在全息描述中，这对应于沿 AdS 边界为 $S_{\mathrm{bulk}}$ 中出现的规范场选取边界条件。

我们讨论两个例子：

\begin{examplebox}{例 6.3：$\mathrm{AdS}_5 \times S^5$：4d $\mathcal{N}=4$ SYM}
考虑 10d IIB 超引力中著名的 $\mathrm{AdS}_5 \times S^5$ 构造，它对偶于规范代数为 $\mathfrak{su}(N)$ 的 4d $\mathcal{N}=4$ 纯超 Yang--Mills 理论。从 10d IIB 超引力解出发，把高阶形式场在背景解附近展开。把这些涨落代入 10d 超引力的运动方程与 Bianchi 恒等式，就得到一组方程。由于我们关心边界处的主导贡献，把它们截断到单导数项，即大距离处的主导项。由此可以反推出重现这些运动方程的拓扑作用量，得到 [217]
\begin{equation}
  S_{\mathrm{top}} = N \int_{\mathrm{AdS}_5} b_2 \wedge \dd c_2\,. \tag{6.31}
\end{equation}
关于该步骤的更多细节见 [215]。这里 $b_2,c_2$ 是连续的体 2-形式规范场，它们来自把 IIB 超引力的 NSNS 与 RR 2-形式规范场在固定背景解附近展开。

作用量 $S_{\mathrm{top}}$ 描述规范代数为 $\mathfrak{su}(N)$ 的 $\mathcal{N}=4$ 理论的 1-形式对称性的 SymTFT。下面考虑边界条件的可能选择。

\para{(a) 边界条件 1。}
边界条件的一种简单选择是
\begin{equation}
  b_2\ \text{Dirichlet}\,,\qquad c_2\ \text{Neumann}\,. \tag{6.32}
\end{equation}

% [原文 p.136]
作用量 (6.31) 的体运动方程是
\begin{equation}
  N b_2 = 0\,,\qquad N c_2 = 0\,. \tag{6.33}
\end{equation}
与这些运动方程相容，意味着这些边界条件迫使 $c_2$ 在边界内以 $\Z_N$ 为模变化。相应地，边界上这个 $\Z_N$-值规范场的 Wilson 面产生一个 1-形式对称性，也就是说我们选择了所得 4d 理论的、带有 1-形式对称性
\begin{equation}
  G^{(1)} = \Z_N \tag{6.34}
\end{equation}
的一个全局形式。我们断言：这些边界条件对应于选择规范群为
\begin{equation}
  \mathcal{G} = \SU(N) \tag{6.35}
\end{equation}
的全局形式。这一点可以在全息上通过考察 7-形式通量的 Bianchi 恒等式看出，它充当 D5-膜的源。特别地，在紧空间 $S^5$ 上积分有
\begin{equation}
  \int_{S^5} \dd F_7 = \int_{S^5} F_5 \wedge H_3 = N H_3\,. \tag{6.36}
\end{equation}
这意味着 $N$ 条基本弦（它们耦合到 $b_2$）可以通过终止于包裹 $S^5$ 的 D5-膜而被屏蔽，而后者正是 [218] 的重子顶点。基本弦成为 4d 规范理论中的 Wilson 线，于是我们得知 (6.34) 纯粹是电 1-形式对称性，从而 (6.35) 得到说明。

\para{(b) 边界条件 2。}
现在交换 $b_2$ 与 $c_2$ 的 Neumann 与 Dirichlet 性质，施加边界条件
\begin{equation}
  c_2\ \text{Dirichlet}\,,\qquad b_2\ \text{Neumann}\,. \tag{6.37}
\end{equation}
由与上面类似的逻辑可知，4d 理论应具有
\begin{equation}
  G^{(1)} = \Z_N \tag{6.38}
\end{equation}
的 1-形式对称性。不过这次全局形式是
\begin{equation}
  \mathcal{G} = \PSU(N)\,, \tag{6.39}
\end{equation}
即离散 theta 角平凡的规范群。这同样可以在全息上由另一个 7-形式的 Bianchi 恒等式来说明
\begin{equation}
  \int_{S^5} \dd H_7 = -\int_{S^5} F_5 \wedge F_5 = -N F_3\,, \tag{6.40}
\end{equation}
它意味着 $N$ 条 D1 弦（它们耦合到 $c_2$）可以通过终止于包裹 $S^5$ 的 NS5-膜而被屏蔽。由于 D1 弦成为 4d 规范理论中的 't Hooft 线，我们得知 (6.38) 纯粹是磁 1-形式对称性，从而 (6.39) 得到说明。

% [Tutor Note: (6.36) 的右端可用通量量子化直接核对：$\mathrm{AdS}_5\times S^5$ 中
%  $\int_{S^5}F_5=N$（D3-膜荷），故 $\int_{S^5}F_5\wedge H_3=N H_3$，两端同为 3-形式。
%  同理 (6.40) 中 $\int_{S^5}F_5\wedge F_5$ 只留一个 $F_5$ 的 AdS 分量，给出 $-N F_3$。
%  这两个恒等式正是「电 vs 磁」两种边界条件 (6.32)/(6.37) 的全息判据。]

% [原文 p.137]
\para{(c) 更一般的边界条件。}
更一般的边界条件给出规范群的其他可能全局形式，以及可能的离散 theta 角。关于其中一些的讨论见 [215]。
\end{examplebox}

\begin{examplebox}{例 6.4：$\mathrm{AdS}_4 \times S^7/\Z_k$：3d ABJM 理论}
探测 $\C^4/\Z_k$ 奇点的 $N$ 个 M2-膜堆的低能描述，是一个三维 Chern--Simons 物质规范理论，规范群为 $U(N)_k \times U(N)_{-k}$ [219]。大 $N$ 全息对偶是 $\mathrm{AdS}_4 \times S^7/\Z_k$ 上的 11d 超引力。可能的全局变体有很多 [220,221]
\begin{equation}
\begin{aligned}
&U(N)_k \times U(N)_{-k}\,,\quad SU(N)_k \times SU(N)_{-k}\,,\quad (SU(N)_k \times SU(N)_{-k})/\Z_N\\
&(U(N)_k \times U(N)_{-k})/\Z_{m'}\,,\quad (SU(N)_k \times SU(N)_{-k})/\Z_{n'}\,,
\end{aligned}\tag{6.41}
\end{equation}
其中 $m'$ 是 $k$ 的因子，$n'$ 是 $N$ 的因子。

在 [95] 中，作者们用上面讨论的那类分析为所有这些变体给出了全息起源，其中用的是 10d IIA 超引力。他们论证 4d SymTFT 含有一个 BF 项
\begin{equation}
  S_{4d} = \int_{\mathcal{M}_4} B_2 \wedge \left(N \dd A_{D0} + k \dd A_{D4}\right)\,. \tag{6.42}
\end{equation}
这里 $B_2$ 是 NS-NS 2-形式规范场，而 $A$ 场分别是由 D0 与 D4 膜产生的 1-形式规范场。下面我们说明如何用 IIA 超引力推出该项。对 BF 项施加不同的边界条件，就得到具有不同整体对称性与整体规范群的一整套 3d 理论。

现在我们从 10 维 IIA 型超引力的运动方程推出拓扑作用量 (6.42)。$\mathrm{AdS}_4 \times S^7/\Z_k$ 的 M-理论在 $\mathrm{AdS}_4 \times \mathbb{C}P^3$ 上的 IIA 型弦论中有描述。背景通量如下
\begin{equation}
  \int_{\mathbb{C}P^3} \frac{F_6}{2\pi} \sim N\,,\qquad \int_{\mathbb{CP}^1 \subset \mathbb{C}P^3} \frac{F_2}{2\pi} \sim k\,. \tag{6.43}
\end{equation}
我们将遵循 [215] 中勾勒的一般哲学来推出四维有效拓扑作用量。首先把 IIA 型的形式场按内禀空间的同调展开，并选定对偶化框架（即展开 $F_2$ 还是 $F_8$ 等等）。内禀空间 $\mathbb{C}P^3$ 的同调为
\begin{equation}
\begin{aligned}
H^0(\mathbb{C}P^3;\Z) &= H^2(\mathbb{C}P^3;\Z) = H^4(\mathbb{C}P^3;\Z) = H^6(\mathbb{C}P^3;\Z) = \Z\,,\\
H^1(\mathbb{C}P^3;\Z) &= H^3(\mathbb{C}P^3;\Z) = H^5(\mathbb{C}P^3;\Z) = 0\,.
\end{aligned}\tag{6.44}
\end{equation}
我们选择用形式 $\omega_2,\omega_4$ 表示 $H^2$、$H^4$，并作归一化使得
\begin{equation}
  \omega_2 \wedge \omega_4 = \vol_{\mathbb{C}P^3}\,. \tag{6.45}
\end{equation}

% [Tutor Note: (6.43) 的下标写法在原文中不统一：第一个积分写成 $\mathbb{C}P^3$（$\mathbb{C}$ 为黑板体、
%  $P$ 为数学斜体），第二个积分的下限写成 $\mathbb{CP}^1$（$P$ 也是黑板体）。
%  译文照原文保留这一差别，读者不必以为漏了排版。]

% [原文 p.138]
此外我们选择按 $F_2$ 与 $F_6$ 展开，因为对它们有显式的通量量子化条件，而无需借助 Hodge 对偶运算
\begin{equation}
\begin{aligned}
F_2 &= k\omega_2 + g_2\,,\\
F_6 &= N\omega_2 \wedge \omega_4 + f_2 \wedge \omega_4 + f_4 \wedge w_2\,,\\
H_3 &= h_3 + h_1 \wedge w_2\,.
\end{aligned}\tag{6.46}
\end{equation}
$F_4$、$F_8$ 由关系 $F_4 = *F_6$、$F_8 = *F_2$ 得到：
\begin{equation}
\begin{aligned}
F_4 &= N\vol_{\mathrm{AdS}_4} + *_{4}f_2 \wedge \omega_2 + *_{4}f_4 \wedge \omega_4\,,\\
F_8 &= k\vol_{\mathrm{AdS}_4} \wedge \omega_4 + *_{4}g_2 \wedge \omega_2 \wedge \omega_4\,,
\end{aligned}\tag{6.47}
\end{equation}
其中我们用到了恒等式 $*_6\omega_2 = \omega_4$ 与 $*_{10}\vol_{\mathrm{AdS}_4} = \omega_2 \wedge \omega_4$。现在把这些展开代入 IIA 型的运动方程与 Bianchi 恒等式。我们要挑出主导阶的拓扑贡献：凡是不拓扑的项、或含一个以上导数的项，都立刻丢弃。循环考察各运动方程，得到
\begin{equation}
\begin{aligned}
\dd *_{10} F_2 &= H_3 \wedge F_6\\
0 &= N h_3 \wedge \omega_2 \wedge \omega_4\,,
\end{aligned}\tag{6.48}
\end{equation}
以及
\begin{equation}
\begin{aligned}
\dd *_{10} H_3 &= F_2 \wedge F_6\\
0 &= (k f_2 + N g_2) \wedge \omega_2 \wedge \omega_4\,,
\end{aligned}\tag{6.49}
\end{equation}
以及
\begin{equation}
\begin{aligned}
\dd *_{10} F_6 &= H_3 \wedge F_2\\
0 &= k h_3 \wedge \omega_2\,.
\end{aligned}\tag{6.50}
\end{equation}
由 Bianchi 恒等式 $\dd F_6 = H_3 \wedge F_4$、$\dd F_2 = H_3 \wedge F_0$、$\dd H_3 = 0$，我们分离出如下相关贡献（令 $F_0 = 0$）
\begin{equation}
  \dd f_2 = 0\,,\qquad \dd h_3 = 0\,,\qquad \dd g_2 = 0\,, \tag{6.51}
\end{equation}
于是可以定义规范场 $a_1,b_2,c_1$
\begin{equation}
  f_2 = \dd a_1\,,\qquad h_3 = \dd b_2\,,\qquad g_2 = \dd c_1\,. \tag{6.52}
\end{equation}
因此只剩下三个拓扑的、单导数的贡献
\begin{equation}
\begin{aligned}
N \dd b_2 &= 0\,,\\
k \dd b_2 &= 0\,,\\
k \dd a_1 + N \dd c_1 &= 0\,,
\end{aligned}\tag{6.53}
\end{equation}
它们正是 (6.42) 的运动方程，对应如下识别
\begin{equation}
  b_2 \leftrightarrow B_2\,,\qquad c_1 \leftrightarrow A_{D0}\,,\qquad a_1 \leftrightarrow A_{D4}\,. \tag{6.54}
\end{equation}
\end{examplebox}

% [Tutor Note: (6.47) 里 $*_{4}$ 是四维 AdS$_4$ 上的 Hodge 对偶（原文只在此处用下标标出维数，
%  未加定义）。据此可核对次数：$*_4 f_2$ 是 4d 中的 2-形式，故 $*_4f_2\wedge\omega_2$ 为 $2+2=4$ 形式；
%  $*_4 f_4$ 是 0-形式，$*_4 f_4\wedge\omega_4$ 为 4 形式；$*_4 g_2 \wedge \omega_2\wedge\omega_4$ 为
%  $2+6=8$ 形式——与 $F_4$、$F_8$ 的次数吻合，说明 $*_4$ 不可能是十维对偶。]
% [Tutor Note: (6.46) 第二、三行里 $\omega_2$ 被印成了 $w_2$（第三行的 $h_1\wedge w_2$、第二行的
%  $f_4\wedge w_2$），字形与第一行的 $\omega_2$（希腊字母 omega）明显不同，应为原文笔误。
%  判据：(6.45) 已把 $\omega_2$ 定义为 $H^2(\mathbb{C}P^3)$ 的生成元，且 (6.48)--(6.50) 的
%  「$0 = \cdots$」各行一律用 $\omega_2\wedge\omega_4$；若 $w_2$ 是独立的 2-形式，
%  (6.46) 的 $F_6$ 就无法用 CP$^3$ 的同调基展开。译文照原文保留 $w_2$ 以存其真。]
